
\documentclass[11pt,oneside]{book}
\usepackage[T1]{fontenc}
\usepackage[utf8]{inputenc}
\usepackage[margin=1in]{geometry}
\usepackage{hyperref}
\usepackage{listings}
\usepackage{enumitem}

\title{ROS 2 Engineering Curriculum \\ Volume 1: Fundamentals}
\author{Self-Study Edition}
\date{\today}

\begin{document}
\maketitle
\tableofcontents

\chapter{Learning Roadmap}
This volume takes a learner from zero knowledge of ROS 2 to the point where
they can create packages, nodes, publishers, subscribers, services, actions,
and launch files.

\chapter{What is ROS 2?}
\section{Why Robotics Needs Middleware}
Robots consist of many independent subsystems:
sensors, perception, planning, control, and actuators.
ROS 2 provides a common communication framework.

\section{ROS 1 vs ROS 2}
Topics:
\begin{itemize}
\item DDS
\item Discovery
\item Security
\item Real-time support
\item Embedded support
\end{itemize}

\section*{Exercises}
Explain why a robot should be split into multiple nodes.

\chapter{Installing ROS 2}
Ubuntu installation workflow, environment setup,
workspace creation, and verification.

\chapter{ROS 2 Architecture}
Layers:
\begin{enumerate}
\item Application
\item rclcpp/rclpy/rclc
\item rcl
\item rmw
\item DDS
\item Network Transport
\end{enumerate}

Detailed discussion of responsibilities of each layer.

\chapter{Nodes}
Node lifecycle:
\begin{enumerate}
\item Process creation
\item ROS initialization
\item Node creation
\item Communication setup
\item Execution
\item Shutdown
\end{enumerate}

Lab: Create a node in C++ and Python.

\chapter{Topics}
Publisher/subscriber communication model.

Lab:
Create a temperature publisher and subscriber.

Quiz:
Difference between synchronous and asynchronous communication.

\chapter{Messages}
Message definitions, standard interfaces,
custom messages, and serialization basics.

Lab:
Create a custom message package.

\chapter{Services}
Request-response communication.

Lab:
Build a calculator service.

\chapter{Actions}
Goal, feedback, result model.

Lab:
Build a countdown action server.

\chapter{Parameters}
Runtime configuration management.

Lab:
Create configurable nodes.

\chapter{Launch Files}
Python launch system, node orchestration,
parameter injection, remapping.

Lab:
Launch three communicating nodes.

\chapter{Development with C++}
rclcpp fundamentals, node classes,
publishers, subscribers, services,
actions, timers, parameters.

Programming exercises included.

\chapter{Development with Python}
rclpy fundamentals, callback design,
timers, services, actions.

Programming exercises included.

\chapter{Workspace and Build System}
colcon, ament, package.xml,
CMakeLists.txt, dependency management.

Lab:
Build a multi-package workspace.

\chapter{Debugging and Introspection}
Commands:
\begin{verbatim}
ros2 node list
ros2 topic list
ros2 topic echo
ros2 topic info
ros2 doctor
\end{verbatim}

Lab:
Investigate a running ROS graph.

\chapter{Mini Projects}

Project 1:
Robot health monitor.

Project 2:
Sensor dashboard.

Project 3:
Multi-node telemetry system.

\chapter{Final Assessment}

Written questions:
\begin{enumerate}
\item Explain the ROS graph.
\item Explain publishers and subscribers.
\item Compare services and actions.
\item Describe message serialization.
\item Describe package dependencies.
\end{enumerate}

Practical exam:
Create a complete ROS 2 application from scratch.

\end{document}
